
Converting transformers to sub-quadratic architectures enables efficient inference but typically degrades few-shot adaptation capability. This work investigates whether task conditioning integrated during conversion can preserve adaptation. We introduce Task-Conditioned Selective State Space (TC-SSM), which modulates Mamba's state space parameters based on task embeddings discovered via unsupervised clustering of transformer hidden states. Experiments on SuperGLUE classification tasks show that K-means clustering of BERT hidden states achieves cluster purity of 0.74, indicating task-relevant structure. InfoNCE-trained task embeddings yield 29.21\% linear probe accuracy on 6-way task classification (versus 16.67\% random baseline). Low-rank modulation (rank 32) of Mamba's $\Delta$, B, and C matrices produces significant state variance differences across tasks (F=100.35, p<0.0001) with 0.85x inference time relative to vanilla Mamba. In 16-shot evaluation, TC-SSM achieves 51.49\% average accuracy compared to 51.75\% for a transformer baseline, a gap of 0.25\%, reaching 95\% of fine-tuned ceiling in 14 gradient steps. These results suggest that task-conditioned conversion can preserve adaptation capability in sub-quadratic architectures at this scale, though generalization to larger models and task types beyond classification remains untested.

